% Generated from validated Article Draft Result. Do not hand-edit; revise the structured draft instead.
\section{答えるモデルから接続するモデルへ――LaMDA、RETRO、WebGPT}
\noindent\textbf{LaMDA、RETRO、WebGPTは、いずれも『外部情報を使う対話AI』としてまとめるべきではない。LaMDAはdialogue-model framing、RETROはretrieval augmentation、WebGPTはweb browsingとhuman feedbackを組み合わせたgroundingという異なる接続様式を示す。2021年には、parameter内部だけで答えを完結させない設計が複数の方向から現れていた。}\autocite{src-66643bb115ae,src-6ab2302937b2,src-72eada5e8850,src-a3971b3f646f,src-bb1464aabad9}

LaMDAが年次記録で担うのは、対話そのものをmodel designの中心に置くdialogue-model framingである。これはretrieval systemやbrowser toolと同義ではない。会話の文脈でどのような応答を生成するかというinterface-orientedな問題が、一般的なtext completionから切り出されて扱われるようになった点に意味がある。本稿ではGoogle側が示したcapabilityやquality評価を独立検証せず、dialogueが独立した設計対象になったことを記録する。\autocite{src-a3971b3f646f}


RETROは、model parameterを増やすことだけに知識や能力を託すのではなく、retrievalを組み合わせる設計を示した。年次的には、scaleの章で見たdense/sparse architectureとはさらに別に、外部の情報資源へアクセスすることでmodel behaviorを構成する選択肢が前景化したことが重要である。retrievalはparameter-only scalingの代替または補完として扱えるが、後年の一般的なtool-use productそのものへ読み替えるべきではない。\autocite{src-66643bb115ae,src-6ab2302937b2}


WebGPTは、web browsingを介して外部情報へ接続し、citationとhuman feedbackを含むgrounding-orientedな設計を扱う。RETROのretrievalがmodel architecture側の接続として読めるのに対し、WebGPTでは情報探索を伴うinteractionと、その出力を人間側の評価へ接続する構造が前面に出る。両者はともにparameter外の情報を使うが、同じ仕組みではない。\autocite{src-72eada5e8850,src-bb1464aabad9}


三者を並べると、2021年に『接続』という設計が少なくとも三つの意味を持っていたことが見える。dialogueは利用者との継続的な文脈接続、retrievalはmodelと外部corpusの接続、web browsing plus feedbackはmodelと動的な情報探索・人間評価の接続である。これらは後年一つのassistant productに同居し得るが、2021年時点では別々の研究問題として現れていた。その分離を保つことで、統合後のsystemを歴史的に分解して理解できる。\autocite{src-66643bb115ae,src-6ab2302937b2,src-72eada5e8850,src-a3971b3f646f,src-bb1464aabad9}


とりわけretrievalとtool useは区別したい。RETROの存在から、2021年に現在的なtool-using agentが成立していたとは言えない。WebGPTのbrowser interactionも、後年の汎用tool ecosystemをそのまま意味するわけではない。一次資料から言えるのは、model内部のparameterだけに依存せず、外部情報を取得・参照しながら出力を構成する設計が複数の形で技術記録に現れていたことまでである。\autocite{src-66643bb115ae,src-6ab2302937b2,src-72eada5e8850,src-bb1464aabad9}


\begin{claimboundary}
LaMDA、RETRO、WebGPTについて各組織・著者が示した能力、品質、性能、安全性の評価は独立再現したものではない。本稿はdialogue、retrieval、web-grounded interactionという設計境界を一次資料に基づいて整理する。\autocite{src-66643bb115ae,src-6ab2302937b2,src-72eada5e8850,src-a3971b3f646f,src-bb1464aabad9}
\end{claimboundary}
